
La carga repartida sobre $AB$ equivale a:
\[
Q=qL=4(3)=\qty{12}{\kilo\newton}
\]
aplicada a $\qty{1.5}{\metre}$ de $A$.

Equilibrio vertical de $AB$:
\[
R_A-12-3=0
\]
y por tanto:
\[
\boxed{R_A=\qty{15}{\kilo\newton}}.
\]

Momentos sobre $A$:
\[
M_A-12(1.5)-3(3)=0,
\]
luego:
\[
\boxed{M_A=\qty{27}{\kilo\newton\metre}}.
\]

La condición esencial es:
\[
\boxed{M_B=0}
\]
por existir una rótula interna.

\begin{tip}
Cuando aparezca una rótula interna, muchas veces conviene \textbf{separar la estructura en dos}. La ecuación $M_{\text{rótula}}=0$ es información adicional muy valiosa.
\end{tip}
\end{solucion}

\begin{ejercicio}[10: Ala idealizada como voladizo]
Una semiala se idealiza como un voladizo de $L=\qty{8}{\metre}$. La sustentación es triangular, máxima en la raíz:
\[
q_L(0)=\qty{15}{\kilo\newton\per\metre}
\]
y nula en la punta.

Calcular el cortante y el momento transmitidos en la raíz.
\end{ejercicio}

\begin{solucion}
La fuerza aerodinámica total es:
\[
L_{\mathrm{tot}}=\frac12q_0L
=\frac12(15)(8)
=\boxed{\qty{60}{\kilo\newton}}.
\]
El centroide está a:
\[
\frac{L}{3}=\frac83=\qty{2.667}{\metre}
\]
desde la raíz.

El momento aerodinámico respecto a la raíz es:
\[
60\frac83=\boxed{\qty{160}{\kilo\newton\metre}}.
\]

Por tanto, en magnitud, la unión ala--fuselaje debe transmitir:
\[
\boxed{|V_{\text{raíz}}|=\qty{60}{\kilo\newton}}
\]
y
\[
\boxed{|M_{\text{raíz}}|=\qty{160}{\kilo\newton\metre}}.
\]

\begin{idea}
Esto explica por qué la zona de raíz del ala es estructuralmente crítica: al acercarnos al fuselaje se ha acumulado prácticamente toda la carga aerodinámica de la semiala.
\end{idea}
\end{solucion}

% ============================================================
\safechapter{Problemas para practicar}
% ============================================================

Resuelve primero estos ejercicios sin mirar los resultados.

\Needspace{5\baselineskip}\section*{Problema 1}
Viga simplemente apoyada:
\[
L=\qty{5}{\metre},\qquad q=\qty{4}{\kilo\newton\per\metre}.
\]
Obtener reacciones, $V(x)$, $M(x)$ y $M_{\max}$.

\subsection*{Resultado}
\[
R_A=R_B=\qty{10}{\kilo\newton}
\]
\[
V(x)=10-4x
\]
\[
M(x)=10x-2x^2
\]
\[
M_{\max}=\boxed{\qty{12.5}{\kilo\newton\metre}}
\quad\text{en}\quad
x=\qty{2.5}{\metre}.
\]

\Needspace{5\baselineskip}\section*{Problema 2}
Voladizo:
\[
L=\qty{2.5}{\metre}
\]
con $P=\qty{8}{\kilo\newton}$ descendente en el extremo libre.

\subsection*{Resultado}
\[
R_A=\qty{8}{\kilo\newton}
\]
\[
M_A^{\text{reacción}}=\qty{20}{\kilo\newton\metre}
